\chapter{Funzione di Rete}

\begin{figure}[h]
    \centering
    \includegraphics[scale = 0.5]{Esempio di un circuito con una funzione di rete.png}
\end{figure}  

\newpage 

\section{Funzioni di Rete}

\footnote{ 
    Slide del prof | 6 - Funzioni di Rete | pag 2\\
    Appunti | 2025-10-15 | pag 3  
}

\begin{tcolorbox}
    Come il caso delle funzioni di rete nel dominio s viste ad Elettrotecnica, 
    adesso analizzeremo, in questo corso di DSP, le funzioni di rete nel dominio z. \newline 

    Per avere un ripasso al volo: \newline 

    Funzione di trasferimento - Wikipedia \\
    \url{https://it.wikipedia.org/wiki/Funzione_di_trasferimento}
\end{tcolorbox} 

Per un circuito, o sistema, a tempo discreto descritto dall'equazione alle differenze finite: 

{
    \Large 
    \begin{equation}
        \sum_{k = 0}^{N}
        a_k \cdot y \left[ n - k\right]
        = 
        \sum_{k = 0}^{M}
        b_k \cdot x \left[ n - k\right]
    \end{equation}
}

dove: 

\begin{itemize}
    \item $x \left[ n\right]$ è una sequenza di ingresso al sistema LTI a tempo discreto 
    \item $y \left[ n\right]$ è una sequenza di uscita al sistema LTI a tempo discreto 
    \item $a_k$ e $b_k$ sono coefficienti reali
\end{itemize}

Applicando la trasformata-z al sistema, avremo che: 

{
    \Large 
    \begin{equation}
        \begin{split}
        \sum_{k = 0}^{N}
        a_k \cdot y \left[ n - k\right]
        &= 
        \sum_{k = 0}^{M}
        b_k \cdot x \left[ n - k\right]
        \\
        &\updownarrow
        \\
        \sum_{k = 0}^{N}
        a_k \cdot z^{-k} \cdot Y \left( z\right)
        &= 
        \sum_{k = 0}^{M}
        b_k \cdot z^{-k} \cdot X \left( z\right)
        \end{split}
    \end{equation}
}

dove: 

\begin{itemize}
    \item $X \left( z\right)$ è la trasformata-z della sequenza in ingresso al sistema LTI a tempo discreto $x \left[ n\right]$
    \item $Y \left( z\right)$ è la trasformata-z della sequenza in uscita al sistema LTI a tempo discreto $y \left[ n\right]$
\end{itemize} 

\newpage 

Visto che i termini $Y \left( z\right)$ e $X \left( z\right)$ sono costanti nella sommatoria nel dominio z (a differenza del dominio n), 
possiamo svolgere queste sostituzioni (step-by-step):

{
    \Large 
    \begin{equation}
        \begin{split}
        \sum_{k = 0}^{N}
        a_k \cdot z^{-k} \cdot Y \left( z\right)
        &= 
        \sum_{k = 0}^{M}
        b_k \cdot z^{-k} \cdot X \left( z\right)
        \\
        Y \left( z\right)
        &=
        \dfrac{\sum_{k = 0}^{M}
        b_k \cdot z^{-k}
        }{\sum_{k = 0}^{N}
        a_k \cdot z^{-k}}
        \cdot X \left( z\right)
        \end{split}
    \end{equation}
}

dove: 

{
    \Large 
    \begin{equation}
        H \left(z\right)
        =
        \dfrac{\sum_{k = 0}^{M}
        b_k \cdot z^{-k}
        }{\sum_{k = 0}^{N}
        a_k \cdot z^{-k}} 
    \end{equation}

}

prende il nome di funzione di rete (o in inglese System Function) perchè descrive il legame ingresso-uscita da una funzione razionale. \newline 

\newpage 

\subsection{La funzioni di rete nei vari domini}

\footnote{ 
    Slide del prof | 6 - Funzioni di Rete | pag 3\\
    Appunti | 2025-10-15 | pag 3  
}

Come dimostrato precedentemente, la funzione di rete $H \left(z\right)$ può essere ancora calcolata, 
sapendo la trasformata-z dell'uscita del sistema LTI a tempo discreto $Y \left(z\right)$ 
e  la trasformata-z dell'ingresso del sistema LTI a tempo discreto $X \left(z\right)$ anche come: 

{
    \Large 
    \begin{equation}
        H \left( z\right)
        = 
        \dfrac{Y \left(z\right)}{X \left(z\right)}
        = 
        \dfrac{\sum_{k = 0}^{M}
        b_k \cdot z^{-k}
        }{\sum_{k = 0}^{N}
        a_k \cdot z^{-k}} 
    \end{equation}
}

$H \left(z\right)$ è anche la trasformata-z della risposta impulsiva di un sistema LTI a tempo discreto $h \left[ n\right]$. \newline 

In formule: 

{
    \Large 
    \begin{equation}
        H \left( z\right)
        = 
        Z \left\{ h \left[ n\right]\right\}
    \end{equation}
}

Un circuito (o sistema) a tempo discreto, lineare e stazionario (cioè LTI) è caratterizzato da un legame ingresso-uscita, 
che può essere rappresentato in tre domini: \newline 

\textbf{Dominio - Tempo} \newline 

{
    \Large 
    \begin{equation}
        y \left[ n\right]
        = 
        h \left[ n\right]
        * 
        x \left[ n\right]
    \end{equation}
}

dove: 

\begin{itemize}
    \item $h \left[ n\right]$ è la risposta impulsiva del sistema LTI 
    \item l'uscita del sistema LTI nel dominio n, cioè $y \left[ n\right]$, è una convoluzione (che è un processo matematico difficile da computazionare) tra la risposta impulsiva $h \left[ n\right]$ e l'ingresso $x \left[ n\right]$ \newline
\end{itemize} 

\textbf{Dominio - Frequenza} \newline 

\begin{tcolorbox}
    Tecnicamente dovrebbe essere il dominio delle pulsazione, ma sappiamo bene che tra pulsazioni e frequenza c'è il coefficiente $2 \pi$ che li differisce. \newline 

    Come dice Chiaraluce a lezione di Telecomunicazioni: "Tanto è solo un coefficiente"
\end{tcolorbox}

{
    \Large 
    \begin{equation}
        Y \left( e^{\jmath \omega}\right)
        = 
        H \left( e^{\jmath \omega}\right) 
        \cdot 
        X \left( e^{\jmath \omega}\right) 
    \end{equation}
} 

dove: 

\begin{itemize}
    \item $H \left( e^{\jmath \omega}\right)$ è la DTFT di $h \left[ n\right]$
    \item l'uscita del sistema LTI nel dominio delle frequenze, cioè $Y \left( e^{\jmath \omega}\right)$, è un prodotto algebrico (che è un processo matematico facile da computazionare) tra la risposta in frequenza $H \left( e^{\jmath \omega}\right)$ e l'ingresso $X \left( e^{\jmath \omega}\right)$ \newline
\end{itemize}

\newpage 

\textbf{Dominio - Dominio-z} \newline 

{
    \Large 
    \begin{equation}
        Y \left( z\right)
        = 
        H \left( z\right) 
        \cdot 
        X \left( z\right) 
    \end{equation}
} 

dove: 

\begin{itemize}
    \item $H \left( z \right)$ è la trasformata-z di $h \left[ n\right]$
    \item anche qui, come nel dominio delle frequenze, l'uscita del sistema LTI nel dominio delle z, cioè $Y \left( z \right)$, è un prodotto algebrico (che è un processo matematico facile da computazionare) tra la risposta nel dominio z $H \left( z\right)$ e l'ingresso $X \left( z\right)$ \newline
\end{itemize} 

\begin{tcolorbox}
Ricordiamo al volo che per passare dal dominio delle frequenza al dominio z è necessario definire una ROC, 
importante soprattutto per l'anti-trasformata-z (vedi \ref{sec: Proprietà della ROC} \nameref{sec: Proprietà della ROC} pag \pageref{sec: Proprietà della ROC})   
\end{tcolorbox}

\newpage 

\section{System Function}

\footnote{ 
    Slide del prof | 6 - Funzioni di Rete | pag 5\\
    Appunti | 2025-10-15 | pag 3 - 4  
}

Per sistemi LTI: 

\begin{figure}[h]
    \centering
    \includegraphics[scale = 0.2]{Sistea LTI a tempo discreto.png}
\end{figure}  

con risposta impulsiva $h \left[ n\right]$ ed ingressi Z-trasformabili, vale: 

{
    \Large 
    \begin{equation}
        Y \left( z \right)
        = 
        H \left(z\right)
        \cdot 
        X \left( z\right)
        \quad 
        z \in ROC 
    \end{equation}
}

Se esistono le trasformate di Fourier (quindi si visualizza il sistema LTI non più a regime transitorio, bensì nel regime permanente), 
cioè che è ammissibile scrivere: 

{
    \Large 
    \begin{equation}
        z = e^{\jmath \omega}
    \end{equation}
}

allora possiamo scrivere che valgono le seguenti relazioni: 

{
    \Large 
    \begin{equation}
        \begin{split}
            \abs{Y \left(e^{\jmath \omega}\right)}
            &= 
            \abs{H \left(e^{\jmath \omega}\right)}
            \cdot 
            \abs{X \left(e^{\jmath \omega}\right)}
            \\
            &\quad
            \\
            \sphericalangle Y \left(e^{\jmath \omega}\right)
            &= 
            \sphericalangle H \left(e^{\jmath \omega}\right)
            +
            \sphericalangle X \left(e^{\jmath \omega}\right) 
        \end{split}
    \end{equation}
}

cioè che: 

\begin{itemize}
    \item per quanto riguarda i moduli, il modulo di $Y \left(e^{\jmath \omega}\right)$ è il prodotto tra il modulo di $H \left(e^{\jmath \omega}\right)$ e il modulo di $X \left(e^{\jmath \omega}\right)$ 
    \item invece, per quanto riguarda le fasi, la fase di $Y \left(e^{\jmath \omega}\right)$ è la somma tra la fase di $H \left(e^{\jmath \omega}\right)$ e la fase di $X \left(e^{\jmath \omega}\right)$ 
\end{itemize}

Grazie a queste relazioni, possiamo scrivere che, per quanto riguarda il filtro LTI in Fourier, 
la sua funzione di rete $H \left( e^{\jmath \omega}\right)$ in Fourier è: 

{
    \Large 
    \begin{equation}
        H \left( e^{\jmath \omega}\right)
        = 
        \abs{H \left( e^{\jmath \omega}\right)}
        \cdot 
        e^{\jmath \sphericalangle H \left( e^{\jmath \omega}\right)}
    \end{equation}
}

dove: 

\begin{itemize}
    \item $\abs{H \left( e^{\jmath \omega}\right)}$ è la risposta in ampiezza del filtro LTI 
    \item $e^{\jmath \sphericalangle H \left( e^{\jmath \omega}\right)}$ è la risposta in fase del filtro LTI
\end{itemize}

Sia $\abs{H \left( e^{\jmath \omega}\right)}$ che $e^{\jmath \sphericalangle H \left( e^{\jmath \omega}\right)}$ sono funzioni continue nel dominio $\omega$. \newline 

\newpage 

\subsection{Magnitude}

\footnote{ 
    Slide del prof | 6 - Funzioni di Rete | pag 6\\
    Appunti | 2025-10-15 | pag 4  
}

Nei filtri LTI, spesso si considerano i segnali random che, essendo ciclostazionari, 
non analizziamo il loro modulo (in inglese Magnitude) in Fourier, bensì la loro varianza. \newline 

\begin{tcolorbox}[colback=green!5!white,colframe=green!75!black]
TO-DO: Inserire il capitolo di Telecomunicazioni in cui si dimostra ciò
\end{tcolorbox} 

In particolare, data la trasformata di Fouerier della risposta impulsiva $h \left[ n\right]$, 
cioè $H \left( e^{\jmath \omega}\right)$, consideriamo come spettro dei moduli la sua potenza, cioè: 

{
    \Large 
    \begin{equation}
        \abs{H \left( e^{\jmath \omega}\right)}^{2}
    \end{equation}
}

che prende il nome di "magnitude-squared". \newline 

\begin{tcolorbox}[colback=red!5!white,colframe=red!75!black] \label{eq: 2}
$\abs{H \left( e^{\jmath \omega}\right)}^{2}$, oltre in inglese a "magnitude-squared", quale è il suo equivalente, ha un nome in italiano? \newline 

Spettro dei moduli della potenza? Gli si può assegnare un nome o una variabile? 
\end{tcolorbox}

Generalmente di $\abs{H \left( e^{\jmath \omega}\right)}^{2}$ si calcola il suo valore in decibel, o dB, come (step-by-step): 

{
    \Large 
    \begin{equation}
        \begin{split}
            10 \cdot \log_{10} \left( \abs{H \left( e^{\jmath \omega}\right)}^{2}\right)
            &= 
            2 \cdot 10 \cdot \log_{10} \left( \abs{H \left( e^{\jmath \omega}\right)}\right)
            \\
            &= 
            20 \cdot \log_{10} \left( \abs{H \left( e^{\jmath \omega}\right)}\right)
        \end{split}
    \end{equation}
}

$20 \cdot \log_{10} \left( \abs{H \left( e^{\jmath \omega}\right)}\right)$ prende il nome, in inglese, di "log magnitude". \newline 

\begin{tcolorbox}[colback=red!5!white,colframe=red!75!black] \label{eq: 3}
Stessa domanda di prima, ma riferita a $20 \cdot \log_{10} \left( \abs{H \left( e^{\jmath \omega}\right)}\right)$
\end{tcolorbox}

Consideriamo il "log magnitude" un guadagno in dB se è positivo, cioè: 

{
    \Large 
    \begin{equation}
        20 \cdot \log_{10} \left( \abs{H \left( e^{\jmath \omega}\right)}\right) > 0
    \end{equation}
}

Consideriamo il "log magnitude" un'attenuazione in dB se è negativo, cioè: 

{
    \Large 
    \begin{equation}
        20 \cdot \log_{10} \left( \abs{H \left( e^{\jmath \omega}\right)}\right) < 0
    \end{equation}
}

\begin{tcolorbox}
    Questo accade perchè, ricordando da Teoria dei segnali, se il dB è positivo significa che si moltiplica l'argomento del logaritmo per un numero maggiore di 1. \newline 

    Se il dB è negativo, significa che si moltiplica l'argomento per un numero compreso tra 0 e 1, in pratica si fa una divisione rispetto all'argomento del logaritmo. 
\end{tcolorbox}

\newpage 

Ricordando che, per dei segnali deterministici, avevamo che: 

{
    \Large 
    \begin{equation}
            \abs{Y \left(e^{\jmath \omega}\right)}
            = 
            \abs{H \left(e^{\jmath \omega}\right)}
            \cdot 
            \abs{X \left(e^{\jmath \omega}\right)}
    \end{equation}
}

Per segnali random, avremo che: 

{
    \Large 
    \begin{equation}
        \begin{split}
            \abs{Y \left(e^{\jmath \omega}\right)}
            &= 
            \abs{H \left(e^{\jmath \omega}\right)}
            \cdot 
            \abs{X \left(e^{\jmath \omega}\right)}
            \\
            &\downarrow
            \\
            \abs{Y \left(e^{\jmath \omega}\right)}^{2}
            &= 
            \abs{H \left(e^{\jmath \omega}\right)}^{2}
            \cdot 
            \abs{X \left(e^{\jmath \omega}\right)}^{2}
        \end{split}
    \end{equation}
}

che in dB diventa: 

{
    \Large 
    \begin{equation}
        \begin{split}
            \abs{Y \left(e^{\jmath \omega}\right)}^{2}
            &= 
            \abs{H \left(e^{\jmath \omega}\right)}^{2}
            \cdot 
            \abs{X \left(e^{\jmath \omega}\right)}^{2}
            \\
            &\downarrow
            \\
            20 \cdot \log_{10} \left( \abs{Y \left( e^{\jmath \omega}\right)}\right)
            &=
            20 \cdot \log_{10} \left( \abs{H \left( e^{\jmath \omega}\right)}\right)
            +
            20 \cdot \log_{10} \left( \abs{X \left( e^{\jmath \omega}\right)}\right)
        \end{split}
    \end{equation}
}

Quindi, per segnali random, l'uscita di un filtro LTI in Fourier in dB è una semplice somma. \newline 

In particolare, possiamo graficare lo spettro dei "magnitude-squared" di $H \left(e^{\jmath \omega}\right)$ al variare della pulsazione $\omega$ con il diagramma di Bode. \newline 

\begin{tcolorbox}[colback=red!5!white,colframe=red!75!black] \label{eq: 4}
La figura possiamo considerarla come il diagramma di Bode di $\abs{H \left(e^{\jmath \omega}\right)}^{2}$? \newline 

Se si, inserire riferimenti e link al diagramma di Bode qui.
\end{tcolorbox}

Se si considera solo le pulsazioni positive, possiamo tracciare un esempio di $\abs{H \left(e^{\jmath \omega}\right)}^{2}$ al variare dele pulsazioni $\omega$ (con note): 

\begin{figure}[h]
    \centering
    \includegraphics[scale = 0.35]{Esempio di un filtro passa-basso in db.png}
\end{figure}  

\newpage 

\subsection{Risposta in Fase}

\footnote{ 
    Slide del prof | 6 - Funzioni di Rete | pag 7\\
    Appunti | 2025-10-15 | pag 4  
}

Si definisce la risposta in fase di un sistema LTI in Fourier: 

{
    \Large 
    \begin{equation}
        \arg \left[ H \left( e^{\jmath \omega}\right)\right]
    \end{equation}
}

la quale è una funzione continua in $\omega$ per $\abs{\omega} < \pi$. \newline 

Un esempio di un andamento della risposta in fase di un filtro LTI: 

\begin{figure}[h]
    \centering
    \includegraphics[scale = 0.2]{Esempio dell'andamento della risposta in fase.png}
\end{figure}  

ma, generalmente, si preferisce rappresentarlo nell'intervallo $\left[ - \pi, + \pi\right]$, 
quindi la figura precendente diventa: 

\begin{figure}[h]
    \centering
    \includegraphics[scale = 0.2]{Esempio dell'andamento della risposta in fase nell'intervallo pi.png}
\end{figure}  

Notando la figura, apparentemente esiste una discontinuità, ma nella realtà $arg \left[ H \left( e^{\jmath \omega}\right)\right]$ non lo è. \newline 

Per calcolare la risposta in fase del filtro $arg \left[ H \left( e^{\jmath \omega}\right)\right]$, 
esso sarà uguale a: 

{
    \Large 
    \begin{equation}
        \arg \left[ H \left( e^{\jmath \omega}\right)\right] 
        = 
        \arctan \left( \dfrac{\mathfrak{J} \left[ H \left( e^{\jmath \omega}\right)\right]}{\Re \left[ H \left( e^{\jmath \omega}\right)\right]}\right)
    \end{equation}
}

e va calcolato per tutti e quattro i quadranti. \newline 

A differenza del calcolo del modulo, per la risposta in fase, dati diversi segnali in frequenza, 
non è possibile sommarli. \newline 

Quindi: 

{
    \Large 
    \begin{equation}
        \arg \left[ a + b + c\right]
        \neq 
        \arg[a] 
        + 
        \arg[b] 
        + 
        \arg[c] 
    \end{equation}
}

\newpage 

\subsubsection{Distorsione di Fase}

\footnote{ 
    Slide del prof | 6 - Funzioni di Rete | pag 8\\
    Appunti | 2025-10-15 | pag 4  
}

In molti casi, la risposta in fase di un sistema è considerata una distorsione perchè le varie componenti si sommano con fasi differenti. \newline 

Per questo motivo è necessario distinguere i vari casi della risposta in fase del sistema LTI $\arg \left[ H \left( e^{\jmath \omega}\right)\right]$. \newline 

Di seguito elenchiamo i casi più importanti di $\arg \left[ H \left( e^{\jmath \omega}\right)\right]$. \newline 

\textbf{Fase zero} \newline 

{
    \Large 
    \begin{equation}
        \arg \left[ H \left( e^{\jmath \omega}\right)\right] 
        = 
        0 
        \quad 
        \abs{\omega} < \pi
    \end{equation}
}

In questo caso, non si ha nessuna distorsione perchè il sistema è istantaneo (nella realtà sarà impossibile progettarlo). \newline 

\textbf{Fase lineare} \newline 

{
    \Large 
    \begin{equation}
        \arg \left[ H \left( e^{\jmath \omega}\right)\right] 
        = 
        - \alpha \omega
        \quad 
        \abs{\omega} < \pi, \alpha \text{ costante}
    \end{equation}
}

In questo caso, non si ha nessuna distorsione ed il sistema ha un ritardo di $\alpha$. \newline

\begin{tcolorbox}
    Questo è il caso che vorremmo avere nella maggior parte dei casi perchè 
    $\arg \left[ H \left( e^{\jmath \omega}\right)\right] $ sarà una retta con pendenza negativa (dato da $- \alpha \omega$) lungo le pulsazioni. 
\end{tcolorbox}

Come esempio di risposta in fase lineare: 

{
    \Large 
    \begin{equation}
        \begin{split}
        H \left(e^{\jmath \omega}\right)
        &= 
        2 \cdot e^{- \jmath \omega 3}
        \\
        &\downarrow
        \\
        y \left[ n \right]
        &= 
        2 x \left[ x - 3\right]
        \end{split}
    \end{equation}
}

\begin{tcolorbox}[colback=red!5!white,colframe=red!75!black] \label{eq: 5}
Avremo in uscita questo risultato dato come ingresso una sequenza $x \left[ n\right]$? Quali sono i passaggi logici per arrivare alla $y \left[ n\right]$ data $H \left( e^{\jmath \omega}\right)$ ?
\end{tcolorbox}

\textbf{Fase lineare generalizzata} \newline 

{
    \Large 
    \begin{equation}
        \arg \left[ H \left( e^{\jmath \omega}\right)\right] 
        = 
        - \alpha \omega + \beta
        \quad 
        \abs{\omega} < \pi, \alpha, \beta \text{ costanti}
    \end{equation}
}

In questo caso, non si ha nessuna distorsione,
ma il sistema ha un ritardo di $\alpha$ e costante moltilplicativa di $e^{\jmath \beta}$. \newline

\begin{tcolorbox}
    Come il caso con fase lineare, ma con una costante, che è $\beta$
\end{tcolorbox} 

\newpage 

\textbf{Fase arbitraria} \newline 

{
    \Large 
    \begin{equation}
        \arg \left[ H \left( e^{\jmath \omega}\right)\right] 
        = 
        \text{ non lineare}
    \end{equation}
}

In questo caso, si ha una distorsione di fase. \newline 

\begin{tcolorbox}
    In questo corso, non vorremmo mai avere una distorsione non lineare.  
\end{tcolorbox}

\newpage 

\section{Ritardo di Gruppo e di Fase}

\footnote{ 
    Slide del prof | 6 - Funzioni di Rete | pag 9 - 10\\
    Appunti | 2025-10-21 | pag 1  
}

